%%%%%%%%%%%%%%%%%%%%%%%%%%%%%%%%%%%%%%%%%%%%%%%%%%%%%%%%%%%%%%%%%%%%%%%%%%%%%%%%%%%%%%%%%%%%%%%%%%%%%%
%
%   Filename    : chapter_2.tex 
%
%   Description : This file will contain your review of related works.
%                 
%%%%%%%%%%%%%%%%%%%%%%%%%%%%%%%%%%%%%%%%%%%%%%%%%%%%%%%%%%%%%%%%%%%%%%%%%%%%%%%%%%%%%%%%%%%%%%%%%%%%%%

\chapter{Related Works}
\label{sec:relatedworks}

\textbf{2.1} \textbf{Chatbots Developed to Provide Support}

Conversational agents, or chatbots, have gained popularity across multiple platforms and fields, including mental health, education, and customer service. These AI-driven systems use natural language processing (NLP) to facilitate human-like interactions, providing real-time assistance to their users. Research has shown that chatbots can be a tool to lower communication barriers, particularly for individuals who experience social isolation, and awkwardness, or have difficulty accessing human support systems. Given their role in increasing role in various domains, it becomes critical to examine their applications in mental health and social communication, as these areas highlight the potential of chatbots in addressing emotional and interpersonal challenges. In this section, various chatbot applications designed to provide support, particularly in mental health and social communication, are discussed.


\textbf{2.1.1} \textbf{Distinguishing Social Communication Support and Mental Health Support}

Chatbots have proven to be valuable tools for supporting both social communication and mental health, yet they differ in their design and objectives. Social communication support chatbots primarily aim to enhance interpersonal skills by simulating real-life conversations, offering structured feedback, and providing role playing exercises. Research has shown that social communication difficulties, particularly in childhood and adolescence, are strongly associated with long-term mental health issues, including anxiety, depression, and behavioral problems (Dall et al., 2022). Additionally, children and young adults who struggle with social communication often experience challenges in forming meaningful social relationships, which can lead to isolation and emotional distress. Early interventions aimed at improving social communication skills can mitigate the risk of developing mental health disorders later in life.

On the other hand, mental health support is designed to provide emotional and psychological assistance, targeting conditions such as depression, anxiety, and stress. Unlike social communication support, which focus on enhancing interpersonal interactions, mental health support addresses the internal emotional struggles an individual faces. Digital platforms, including computer-mediated communication (CMC) tools, play a growing role in mental health support by offering accessible interventions. However the effectiveness of these tools for mental health support depends on various factors, including the type of communication used and the context in which it occurs. High et al. (2022) suggest that online interactions can help individuals express emotions more freely, seek support from peers, and gain access to professional help. At the same time, over-reliance on digital communication for emotional support may be a contributing factor to social withdrawal and reduced face-to-face interactions, which are critical for psychological well-being. 


\textbf{2.1.2 Chatbots Developed to Provide Support for Mental Health}

Woebot is a conversational agent designed to provide mental health support through guided self-help strategies based on cognitive behavioral therapy (CBT). By engaging users through daily check-ins, mood tracking, and psychoeducationallessons, Woebot adapts to the user’s emotional state and provides tailored interactions. Through this, Fitzpatrick et al. (2017) found that Woebot effectively reduced depressive symptoms in young adults, as measured by the Patient Health Questionnaire-9 (PHQ-9). Despite its strengths, such as high engagement levels and accessibility, its limitations were also noted, including repetitive responses and a lack of deep conversational understanding, which could affect long-term user satisfaction.

Beyond structured CBT interventions, other chatbots explore more interactive and creative methods to support emotional well-being. Eren et al. (2020) introduced Therapist Vibe, a chatbot that encourages children to recognize and express their emotions through storytelling. By prompting children to describe emotional experiences and reflect on possible solutions, Therapist Vibe fosters emotional intelligence in a safe, non-judgmental space. The study found that children using the chatbot were more willing to discuss their emotions and gained insights into emotional management. This highlights how chatbots can serve not only as mental health support tools but also as educational instruments for emotional development, amongst children in this case.

Additionally,  recent studies have expanded on the acceptability and potential of chatbots for mental health support, especially among young adults. Koulouri, Macredi, and Olakitan 92022) explored the acceptability of chatbots in mental healthcare for young adults, identifying them as a feasible solution. Their study found that young adults view chatbots as a potential tool for mental health support, appreciating their convenience and accessibility. However, the research also highlighted challenges such as stigma and the need for user-centered design. The study emphasized the importance of the chatbot’s development process, despite it being viewed as an acceptable solution. It should involve careful attention to user needs and technological limitations, including making sure that chatbots offer personalized and non-judgmental interactions. This finding aligns with those of Haque and Rubya (2023), who noted that chatbots could foster a judgment-free environment, allowing users to share sensitive information. In their study, they examined various commercially available mental health chatbots by analyzing over 6,000 user reviews from both the Google Play Store and Apple App Store. The study found that while users positively perceived chatbots’ personalized, humanlike interactions, several challenges were noted. For one, improper responses and assumptions about the personalities of users led to diminished user engagement. Additionally, while the always-available nature of chatbots was appreciated, users often became overly reliant on them, sometimes preferring them over interactions with friends or family. Furthermore, despite the 24/7 availability of crisis care features, chatbots still lacked the capability to properly identify and respond to crisis situations. It was also noted that some chatbots were unable to understand the use of language associated with a mental health crisis. On the other hand, the study also emphasized the judgment-free environment created by chatbots, which helped users feel more comfortable sharing sensitive information. These findings highlight the potential of chatbots to provide support in situations where real-world human interaction is difficult for users to achieve, though it also emphasizes that excessive reliance on chatbots could lead to social isolation or insufficient crisis care.


\textbf{2.1.3 Chatbots Developed to Provide Support for Social Communication}

In addition to their role in mental health, chatbots have been employed to improve users’ social communication skills, particularly in professional and educational settings. One such example is an experiment conducted for the security company “Workrate,” which faced challenges in training its recruit security guards due to expensive, time-consuming, and inconsistent training methods. De Bever et al. (2019) developed two versions of a chatbot: one that accepts voice and text input and another that only accepts text. With this, participants were divided into three groups: one using the voice and text chatbot, another using the text-only chatbot, and a control group undergoing traditional learning methods. While no significant differences in exam scores were found between the groups, qualitative data showed that participants appreciated the chatbot’s engaging and interactive nature. This study emphasizes the potential of chatbots as supplemental training tools, particularly in fields where interpersonal communication is involved.

Further advancing the concept of chatbot-assisted social skill development, researchers introduced the AI Partner-AI Mentor (APAM) framework. This system integrates an AI Partner, which roleplays human interactions to provide experiential training, and an AI Mentor, which offers personalized feedback and factual guidance. The APAM framework is designed to create an effective learning environment by following four key steps: understanding the social processes underlying the target skill, designing an AI Partner to simulate realistic conversations, developing an AI Mentor to provide tailored feedback, and integrating both agents into a simulated environment for iterative learning. While previous studies have explored the use of large language models (LLMs) in social skills training, APAM remains in the conceptual stage, with empirical testing yet to be conducted. The study highlights the need for LLM-based training chatbots to be unbiased, ethical, and context-aware while maintaining consistent, plausible, and instructable personalities. Additionally, concerns regarding job displacement highlight the importance of a gradual, risk-aware deployment of APAM, emphasizing human-centered design  and continuous evaluation to maximize its effectiveness.

\hl{Currently, a notable AI employed to assist users' communication skills is Duolingo's roleplay feature wherein users are given a situation and are asked to converse with a chatbot. They would then be assessed based on how they communicated and given feedback based on their answers with this chatbot (Team D, 2023). Connecting this to the previous paper, we can see that this is an implementation of the AI Partner, thereby it will be an example of how social communication is defined not by mechanical execution (e.g., syntax or vocabulary use), but rather by the user's ability to select contextually appropriate responses.}

Recent studies have also highlighted how social chatbots can support individuals with social deficits, particularly those with autism spectrum disorder (ASD) or social anxiety. Franze et al. (2023) discuss how AI-driven chatbots such as ChatGPT, can provide a safe and nonjudgmental space for individuals with ASD to practice social interactions. These chatbots give the users the space to rehearse conversations, develop confidence, and receive real-time feedback without the fear of negative social consequences. However, the same study discusses the potential risks associated with chatbot use, including the reduced engagement in real-world interactions, user dependency on the chatbot, and reinforcement of dysfunctional social behaviors. Some chatbots, designed to be highly responsive to user input, may inadvertently encourage conversational habits that are unfavorable to real-world social communication, such as being in control of the entire dialogue, skipping responses, or fixating on a limited range of conversational interests.

\textbf{2.2} \textbf{Other Technologies for Social Skills Training}

Other technologies were also leveraged to help improve social communication skills. Hoque et al. (2014) developed “My Automated Conversation Coach” or MACH, a fully automated nonverbal sensing system designed to improve social interaction skills. MACH interprets and provides real-time feedback to an individual’s nonverbal behavior. It tackles challenges in regards to the technology’s ability to recognize nonverbal cues, applications for social skills training, and potential future developments for automated feedback systems. A simulated job interview was done to assess and gauge the participants’ social skills, comparing the skills of those who used MACH, those who only reviewed MACH recordings, and those who didn’t use MACH. The use of MACH with its automated feedback resulted in a statistically significant improvement in social skills.

DeRosier et al. (2019) makes use of a game developed by 3C Institute called “Hall of Heroes” which provides social skills training through a virtual story set in a superhero-themed world.  Participants were split into two groups, one that played the game, and one that did not. Results showed a significant change in the experimental group’s interpersonal skills, with improvement in the ability to relate to others, accept affection, and express emotion. Additionally, there was also a significant decline in feelings of anxiety, depression, and hopelessness in the experimental group compared to the control group. Evaluations done by both the children and the parents suggest positive reception towards using the game for social skills training, indicating that video games can serve as an alternative and favorable approach towards SST.

Recent developments in the virtual world also led to opportunities to make use of it in aiding with an individual’s social communication skills. A pilot study was conducted by Kolk et al. (2022) that involves a combination of multitouch-multiuser tabletops (MMT) and PowerVR, a virtual reality platform, and their use for training social communication skills for children with neurological disorders. The training sessions were split into one hour on the tabletop, and 30 minutes using VR technology once a week for ten weeks, with two therapists guiding them. Results showed improvements in the children’s ability to start conversations, as well as improvements in other social aspects such as empathy, conflict resolution, and being able to make friends. Qualitative results also showed positive reception towards the use of VR, with participants saying that its use is interesting, exciting, comfortable, and an innovative environment to learn and practice in various social situations. 

Another technology used for social skills training comes from Lee et al.(2023), who investigates the use of metaverse-based social skills training programs to support children with Autism Spectrum Disorder (ASD) in enhancing their social interaction abilities. Conducted as an open-label, single-center, randomized controlled pilot trial, the research explores how immersive virtual environments can serve as effective tools for social practice. The study highlights the need for personalization in chatbot-driven interactions and evaluates how participants respond to AI-generated feedback. By reviewing the opinions of participants, the study was able to underscore the importance of chatbot responsiveness and the quality of responses to facilitate effective social skills training. Moreover, the study suggests that future research should focus on refining chatbot content to improve user engagement.

\textbf{2.3} \textbf{Ethical Considerations in AI-Based Social Communication}

The use of artificial intelligence (AI) in social communication has given rise to a number of ethical concerns, particularly in the areas of bias, disinformation, mental health, education, and accessibility. As AI-powered chatbots continually get built and are integrated into society, it is vital to examine the ramifications of their use across different fields and avenues while stressing the necessity of fairness, openness, and human supervision to ensure moral implementation.

The potential for AI-driven chatbots, especially large language models (LLMs) like ChatGPT, to inherently include social prejudices in their responses due to this being present in their training data is a significant ethical worry. According to Zhou et al. (2024), biased AI answers can perpetuate discrimination and false information thus calling for tactics like dataset audits and fairness-aware model training. Therefore, it is vital that the causalities of AI-decisions are made transparent to all users. This would ensure human-chatbot interactions would avoid bias and foster inclusivity. Moreover, Zhou et al. (2024) advocates for stricter AI ethics policies and frameworks that promote responsible and equitable chatbot interactions in reaction to the findings.

Despite the ease of access provided by chatbots when it comes to accessing support, it is also scrutinized for its dependence on preset scripts and the lack of contextual awareness. Thus, it has problems regarding the provided support’s accuracy, privacy, and possible misinformation (Abd-Alrazag et al., 2021). It is then important to stress that human supervision is required in order to avoid subjecting users to faulty support. Rather than used as a stand-alone tool, LLMs might be best used as a supplementary tool (Abd-Alrazag et al., 2021). Moreover, AI-generated content in educational settings has broader implications for learning experiences, requiring clear ethical guidelines to balance technological benefits with responsible use. 

In addition to this, Saylam et al., (2023) further talks about the challenges regarding AI’s role in education when it comes to academic integrity and critical thinking. The misuse of AI tools may impact their ability to develop their problem-solving skills, thereby increasing their dependency on tools instead of their own academic proficiency. Due to this, it is vital to propose policy frameworks that encourage ethical AI while reaping the benefits when it comes to efficiency and automation of the learning process (Saylam et al., 2023). It is also proposed by Saylam et al., (2023) that a learning environment utilizing and integrating AI could enhance educational outcomes, thus it is important to continually research this for the future. 

As we have assessed in a previous section, the use of AI to support individuals who are part of the spectrum (ASD) may be ethically problematic. Gagan et al. (2023) investigates the potential of AI-based narrative tools to assist in the development of social communication skills in children with Level 1 ASD. The improper implementation of AI to support those that are part of the spectrum, however, may lead to the individual having reduced human empathy, transparency loss, and dependence on users. Gagan et al. (2023) suggest a hybrid approach that overcomes these challenges by enabling AI systems to support human professionals instead of replacing them. This would then allow the implementation and deployment of AI in sensitive applications such as education and mental health.

In totality, the ethical challenges asociated with AI in social communication are related to the importance of transparency, fairness, and human collaboration to ensure that AI does not reinforce bias, compromise privacy, and diminish human empathy. This would then allow for their ethical and responsible deployment in different fields..

{\textbf{2.4} \textbf{Foundational Studies on Social Communication and Mental Health}
}

\textbf{2.4.1 An Experimental Study towards Young Adults: Communication Skills Education}

Guclu (2016) investigates the effectiveness of learning communication skills among young adults, i.e., university students (aged 18 - 24 with an average age of 19.47), and its effects on their empathetic nature and emotional expression. The study illustrates the importance of communication skills in various settings like healthcare, professional settings, and social settings. The results reveal that andragogy-based communication skills training can result in a high level of change in students' empathetic nature and emotional expression. The outcomes reveal that young adults are a suitable target group for social communication training since they are at a developmental phase where the skills are crucial and training can yield measurable gains. The study concludes that communication skills education enhances the functional communication abilities of young adults, which is essential to survive the challenges and demands of life.

\textbf{2.4.2 The Importance of Feeling Awkward: A Dialogical
Narrative Phenomenology of Socially Awkward
Situations}

Clegg (2012) applies dialogical narrative phenomenology to analyze 16 socially awkward experiences. Through analysis, it identifies social awkwardness as characterized by a subjective experience of social or moral transgression, which leads to enhanced social consciousness and a centered focus for social activities. The study participants wished to transcend awkward encounters by staying away from them or addressing them head on. The study concludes that social awkwardness felt and articulated may be an important aspect of social relationships since it can increase sensitivity to social interaction and enable social transformation. This study offers useful information for a chatbot for social communication training. It describes the characteristics of socially awkward situations, defining them as situations that are perceived to be a social or moral transgression and increase social awareness, which can be utilized to make a chatbot recognize and react to the nuances of awkwardness. It also provides insight into awkwardness causes such as off-target speech or social norm violation and action cues such as evasive or anxious behavior,  resolution strategies like avoidance or confrontation. The study also emphasizes the possibility of awkward moments to create social awareness, which can be used by a chatbot to assist users in handling awkwardness and building healthier social interactions. 

\textbf{2.4.3 Awkward but so what: Differences in social trait preferences between autistic and non-autistic adults} 

Dunn et al. (2023) differentiates the perception of autistic and non-autistic adults on their social characteristics, particularly awkwardness. The research involving 24 autistic and 24 non-autistic adults investigated how individuals' social characteristics influence whom they are comfortable with. The findings indicate that autistic adults perceive themselves as more awkward and less competent socially than non-autistic adults but also prefer to be with someone as awkward as them. This disproves the hypothesis that ordinary social behaviors suit everyone. In its conclusion, it states that autistic individuals engage better with the people who also use their type of social approach, highlighting how different social modes must be tolerated and accepted more widely. 

%%%%%%%%%%%%%%%%%%%%%%%%%%%%%%%%%%%%%%%%%%%%%%%%%%%%%%%%%%%%%%%%%%%%%%%%%%%%%%%%%%%%%%%%%%%%%%%%%%%%%%